% Remove the following line and define them in main.tex
%\keywords{M.Sc. degree\spc master’s thesis instructions\spc structure of a master’s thesis\footnote{The keywords should be different from the ones used in the title. Remove this footnote from your thesis.}}

\begin{thesisabstract}{english}
  \abstractmetadata
% Alternatively you can write the abstract here.
% Paragraph breaks and special characters are allowed here, but the content
% should essentially be the same as the abstract in your metadata
Students in the Computer Systems course often struggle to set up their tools and are forced to use simple text prints to find errors in their code. This thesis presents the design and implementation of a unified, containerised debugging platform for the Raspberry Pi RP2040 and RP2350 microcontrollers. The platform utilises Docker containerisation, Virtual Machine technology and a web-based IDE (\textit{code-server}) to provide a scalable solution, ensuring a reproducible education environment that eliminates environmental errors.

A core contribution of this work is the development of the \textit{Pico-Debug-Dashboard}, a custom Visual Studio Code extension that unifies the complexities of \texttt{CMake}, \texttt{OpenOCD}, and \texttt{GDB} into a centralised interface. The platform also supports multiple debugging methods, including physical hardware probes and virtual software probes (\textit{pico-debug}). To evaluate the platform's effectiveness for educational use in the University of Oulu’s Computer Systems course, a comparative analysis was conducted between containerised and virtual-machine-based implementations.

Evaluation was performed using both quantitative metrics—such as deployment latency and flash throughput—and qualitative metrics inspired by the NASA Task Load Index, focusing on beginner-friendliness and ergonomics. The results indicate that the containerised approach significantly reduces environmental cold-start times and provides superior hardware access stability. Among the evaluated debugging probes, the Second Pico hardware implementation was found to be the most suitable for large-scale classroom deployment due to its balance of support, low cost, and dual-core debugging capabilities.
\end{thesisabstract}
